\documentclass[10pt,a4paper]{amsart}
\usepackage{amsmath,amssymb,amsthm,graphicx}
\usepackage{verbatim}
\newenvironment{reference}{\comment}{\endcomment}
%\newenvironment{reference}{}{}
\newenvironment{slogan}{\comment}{\endcomment}
\newenvironment{history}{\comment}{\endcomment}

% For commutative diagrams we use Xy-pic
\usepackage[all]{xy}

% We use 2cell for 2-commutative diagrams.
\xyoption{2cell}
% Theorem environments.
%
\theoremstyle{plain}
\newtheorem{theorem}[subsection]{Theorem}
\newtheorem{proposition}[subsection]{Proposition}
\newtheorem{lemma}[subsection]{Lemma}

\theoremstyle{definition}
\newtheorem{definition}[subsection]{Definition}
\newtheorem{example}[subsection]{Example}
\newtheorem{exercise}[subsection]{Exercise}
\newtheorem{situation}[subsection]{Situation}

\theoremstyle{remark}
\newtheorem{remark}[subsection]{Remark}
\newtheorem{remarks}[subsection]{Remarks}

\numberwithin{equation}{subsection}

% Macros
%
\def\lim{\mathop{\mathrm{lim}}\nolimits}
\def\colim{\mathop{\mathrm{colim}}\nolimits}
\def\Spec{\mathop{\mathrm{Spec}}}
\def\Hom{\mathop{\mathrm{Hom}}\nolimits}
\def\Ext{\mathop{\mathrm{Ext}}\nolimits}
\def\SheafHom{\mathop{\mathcal{H}\!\mathit{om}}\nolimits}
\def\SheafExt{\mathop{\mathcal{E}\!\mathit{xt}}\nolimits}
\def\Sch{\mathit{Sch}}
\def\Mor{\mathop{\mathrm{Mor}}\nolimits}
\def\Ob{\mathop{\mathrm{Ob}}\nolimits}
\def\Sh{\mathop{\mathit{Sh}}\nolimits}
\def\NL{\mathop{N\!L}\nolimits}
\def\CH{\mathop{\mathrm{CH}}\nolimits}
\def\proetale{{pro\text{-}\acute{e}tale}}
\def\etale{{\acute{e}tale}}
\def\QCoh{\mathit{QCoh}}
\def\Ker{\mathop{\mathrm{Ker}}}
\def\Im{\mathop{\mathrm{Im}}}
\def\Coker{\mathop{\mathrm{Coker}}}
\def\Coim{\mathop{\mathrm{Coim}}}

% Boxtimes
%
\DeclareMathSymbol{\boxtimes}{\mathbin}{AMSa}{"02}

%
% Macros for moduli stacks/spaces
%
\def\QCohstack{\mathcal{QC}\!\mathit{oh}}
\def\Cohstack{\mathcal{C}\!\mathit{oh}}
\def\Spacesstack{\mathcal{S}\!\mathit{paces}}
\def\Quotfunctor{\mathrm{Quot}}
\def\Hilbfunctor{\mathrm{Hilb}}
\def\Curvesstack{\mathcal{C}\!\mathit{urves}}
\def\Polarizedstack{\mathcal{P}\!\mathit{olarized}}
\def\Complexesstack{\mathcal{C}\!\mathit{omplexes}}
% \Pic is the operator that assigns to X its picard group, usage \Pic(X)
% \Picardstack_{X/B} denotes the Picard stack of X over B
% \Picardfunctor_{X/B} denotes the Picard functor of X over B
\def\Pic{\mathop{\mathrm{Pic}}\nolimits}
\def\Picardstack{\mathcal{P}\!\mathit{ic}}
\def\Picardfunctor{\mathrm{Pic}}
\def\Deformationcategory{\mathcal{D}\!\mathit{ef}}
\usepackage{cleveref}
\crefname{originallemma}{original lemma}{original lemmas}
\Crefname{originallemma}{Original lemma}{Original lemmas}
\makeatletter
\newlabel{descent-lemma-affine}{{0245}{source}{}{sourcepointer-1}{}}
\newlabel{descent-lemma-affine@cref}{{[originallemma][0][]0245}{[1][0][]source}{}{}{}}
\makeatother
\crefname{originallemma}{original lemma}{original lemmas}
\Crefname{originallemma}{Original lemma}{Original lemmas}
\makeatletter
\newlabel{descent-lemma-descending-property-finite-locally-free}{{02VO}{source}{}{sourcepointer-2}{}}
\newlabel{descent-lemma-descending-property-finite-locally-free@cref}{{[originallemma][0][]02VO}{[1][0][]source}{}{}{}}
\makeatother
\crefname{originallemma}{original lemma}{original lemmas}
\Crefname{originallemma}{Original lemma}{Original lemmas}
\makeatletter
\newlabel{morphisms-lemma-base-change-finite-locally-free}{{02KD}{source}{}{sourcepointer-3}{}}
\newlabel{morphisms-lemma-base-change-finite-locally-free@cref}{{[originallemma][0][]02KD}{[1][0][]source}{}{}{}}
\makeatother
\crefname{originalsection}{original section}{original sections}
\Crefname{originalsection}{Original section}{Original sections}
\makeatletter
\newlabel{section-representable}{{03HA}{source}{}{sourcepointer-4}{}}
\newlabel{section-representable@cref}{{[originalsection][0][]03HA}{[1][0][]source}{}{}{}}
\makeatother
\crefname{originallemma}{original lemma}{original lemmas}
\Crefname{originallemma}{Original lemma}{Original lemmas}
\makeatletter
\newlabel{sites-modules-lemma-local-final-object}{{03DN}{source}{}{sourcepointer-5}{}}
\newlabel{sites-modules-lemma-local-final-object@cref}{{[originallemma][0][]03DN}{[1][0][]source}{}{}{}}
\makeatother
\crefname{originallemma}{original lemma}{original lemmas}
\Crefname{originallemma}{Original lemma}{Original lemmas}
\makeatletter
\newlabel{spaces-lemma-morphism-sheaves-with-P-effective-descent-etale}{{03I2}{source}{}{sourcepointer-6}{}}
\newlabel{spaces-lemma-morphism-sheaves-with-P-effective-descent-etale@cref}{{[originallemma][0][]03I2}{[1][0][]source}{}{}{}}
\makeatother
\crefname{originallemma}{original lemma}{original lemmas}
\Crefname{originallemma}{Original lemma}{Original lemmas}
\makeatletter
\newlabel{spaces-properties-lemma-pushforward-etale-base-change-modules}{{03LX}{source}{}{sourcepointer-7}{}}
\newlabel{spaces-properties-lemma-pushforward-etale-base-change-modules@cref}{{[originallemma][0][]03LX}{[1][0][]source}{}{}{}}
\makeatother
\title{A representable morphism of algebraic spaces is finite locally free exactly when it is affine with finite locally free pushforward of the structure sheaf}
\author{The Stacks Project Authors}
\date{stacks-project:spaces-morphisms.tex:local-snapshot-20261012; GNU FDL 1.2 or later}
\begin{document}
\maketitle
\setcounter{subsection}{1}
\noindent\textit{Editorial note.} All delivered mathematical text is the author's own, copied from the pinned Stacks Project chapter file spaces-morphisms.tex; no statement or proof step was rewritten, reordered or removed.\par\medskip
\noindent\textit{Editorial note.} The statement refers to the representable sense of finite local freeness fixed in Section section-representable of the same chapter; the section label is a navigation pointer and the property is spelled out again in the first line of the proof.\par\medskip
\noindent\textit{Editorial note.} The seven references of the counted spans are recorded as non-load-bearing navigation pointers; the argument that consumes them stays inside the counted proof span.\par\medskip
\noindent\textit{Editorial note.} Preamble: the three native excerpt spans of sources/stacks-spaces-morphisms/preamble.tex (lines 8-13, 14-18 and 154-217). The chapter file itself carries no preamble text.\par\medskip
\noindent\textit{Allowed standard prerequisites (editorial).}\begin{itemize}
\item The base change formula for pushforwards of modules along an etale morphism. Properties of Spaces, Lemma spaces-properties-lemma-pushforward-etale-base-change-modules, used at proof lines 9581-9584 and 9593-9596. Identifies the restriction of f\_*O\_X to V with the pushforward along the base change.
\item Finite local freeness of a module is local on the etale site. Modules on Sites, Lemma sites-modules-lemma-local-final-object, used at proof lines 9585-9588. Concludes finite local freeness of f\_*O\_X from its restriction to a cover.
\item Effective etale descent for morphisms with a sheaf-theoretic property. Spaces, Lemma spaces-lemma-morphism-sheaves-with-P-effective-descent-etale, with Morphisms, Lemma morphisms-lemma-base-change-finite-locally-free and Descent, Lemmas descent-lemma-descending-property-finite-locally-free and descent-lemma-affine, used at proof lines 9598-9604. Descends finite local freeness and affineness from the etale cover back to f.
\item The definition of finite local freeness for representable morphisms. Section section-representable, used at proof lines 9567 and 9573-9578. Fixes the meaning of the property being characterized.
\end{itemize}\medskip
\ifdefined\hypertarget\hypertarget{sourcepointer-1}{}\fi
\noindent\textit{Source reference.} Original lemma 0245: focused-seeds/AGP/sources/stacks-snapshot-20261011/tags.txt, source lines 2766--2766. Navigation pointer into the Stacks Project; the standard fact it supplies is applied inside the counted proof span and the argument resting on it stays there.\par\medskip
\ifdefined\hypertarget\hypertarget{sourcepointer-2}{}\fi
\noindent\textit{Source reference.} Original lemma 02VO: focused-seeds/AGP/sources/stacks-snapshot-20261011/tags.txt, source lines 3775--3775. Navigation pointer into the Stacks Project; the standard fact it supplies is applied inside the counted proof span and the argument resting on it stays there.\par\medskip
\ifdefined\hypertarget\hypertarget{sourcepointer-3}{}\fi
\noindent\textit{Source reference.} Original lemma 02KD: focused-seeds/AGP/sources/stacks-snapshot-20261011/tags.txt, source lines 3368--3368. Navigation pointer into the Stacks Project; the standard fact it supplies is applied inside the counted proof span and the argument resting on it stays there.\par\medskip
\ifdefined\hypertarget\hypertarget{sourcepointer-4}{}\fi
\noindent\textit{Source reference.} Original section 03HA: focused-seeds/AGP/sources/stacks-snapshot-20261011/tags.txt, source lines 4559--4559. Navigation pointer into the Stacks Project; the standard fact it supplies is applied inside the counted proof span and the argument resting on it stays there.\par\medskip
\ifdefined\hypertarget\hypertarget{sourcepointer-5}{}\fi
\noindent\textit{Source reference.} Original lemma 03DN: focused-seeds/AGP/sources/stacks-snapshot-20261011/tags.txt, source lines 4426--4426. Navigation pointer into the Stacks Project; the standard fact it supplies is applied inside the counted proof span and the argument resting on it stays there.\par\medskip
\ifdefined\hypertarget\hypertarget{sourcepointer-6}{}\fi
\noindent\textit{Source reference.} Original lemma 03I2: focused-seeds/AGP/sources/stacks-snapshot-20261011/tags.txt, source lines 4588--4588. Navigation pointer into the Stacks Project; the standard fact it supplies is applied inside the counted proof span and the argument resting on it stays there.\par\medskip
\ifdefined\hypertarget\hypertarget{sourcepointer-7}{}\fi
\noindent\textit{Source reference.} Original lemma 03LX: focused-seeds/AGP/sources/stacks-snapshot-20261011/tags.txt, source lines 4729--4729. Navigation pointer into the Stacks Project; the standard fact it supplies is applied inside the counted proof span and the argument resting on it stays there.\par\medskip
\begin{lemma}
\label{lemma-finite-locally-free-representable}
Let $S$ be a scheme. Let $f : X \to Y$ be a representable morphism
of algebraic spaces over $S$. Then $f$ is finite locally free
(in the sense of Section \ref{section-representable})
if and only if $f$ is affine and the sheaf $f_*\mathcal{O}_X$ is
a finite locally free $\mathcal{O}_Y$-module.
\end{lemma}

\begin{proof}
Assume $f$ is finite locally free (as defined in
Section \ref{section-representable}). This means that
for every morphism $V \to Y$ whose source is a scheme the
base change $f' : V \times_Y X \to V$ is a finite locally free morphism
of schemes. This in turn means (by the definition of a finite locally
free morphism of schemes) that
$f'_*\mathcal{O}_{V \times_Y X}$
is a finite locally free $\mathcal{O}_V$-module. We may choose $V \to Y$
to be surjective and \'etale. By
Properties of Spaces,
Lemma \ref{spaces-properties-lemma-pushforward-etale-base-change-modules}
we conclude the restriction of $f_*\mathcal{O}_X$ to $V$ is
finite locally free. Hence by
Modules on Sites, Lemma \ref{sites-modules-lemma-local-final-object}
applied to the sheaf $f_*\mathcal{O}_X$ on $Y_{spaces, \etale}$
we conclude that $f_*\mathcal{O}_X$ is finite locally free.

\medskip\noindent
Conversely, assume $f$ is affine and that $f_*\mathcal{O}_X$ is a finite
locally free $\mathcal{O}_Y$-module. Let $V$ be a scheme, and let
$V \to Y$ be a surjective \'etale morphism. Again by
Properties of Spaces,
Lemma \ref{spaces-properties-lemma-pushforward-etale-base-change-modules}
we see that $f'_*\mathcal{O}_{V \times_Y X}$ is finite locally free.
Hence $f' : V \times_Y X \to V$ is finite locally free (as it is also affine).
By
Spaces,
Lemma \ref{spaces-lemma-morphism-sheaves-with-P-effective-descent-etale}
we conclude that $f$ is finite locally free (use
Morphisms, Lemma \ref{morphisms-lemma-base-change-finite-locally-free}
Descent, Lemmas \ref{descent-lemma-descending-property-finite-locally-free}
and \ref{descent-lemma-affine}). Thus we win.
\end{proof}

\end{document}
